% !PRINTED 134
\weberpage{134}
%<<CONT
exactly $(p-1)\nu$ elements occur whose degree is divisible by $p$, only the formula

\[ \varepsilon)\quad \nu = j\,r_2 \]

remains to be proved. This results from the following consideration. By $\alpha')$ the number of the elements $U$ is equal to $p\,\nu$. But among these elements $U$ there is certainly one, the unit element, whose degree is not divisible by $p$, and consequently by $\delta)$

\[ j\,r_2\,(p-1) < p\,\nu. \]

On the other hand $n = \mu\,\nu = p\,r_1 r_2 j$, and since $\nu$ is relatively prime to $p\,r_1$, $j\,r_2$ is divisible by $\nu$. Hence $j\,r_2 : \nu$ is a positive integer which is smaller than $p : p-1$ and consequently also smaller than 2, and which can therefore only be equal to 1. Thereby $\varepsilon)$ is proved, and with it the whole theorem.

\begin{center}
\S.\,33.

\gesperrt{Groups of degree $p^{\alpha} q$.}
\end{center}

\gesperrt{Frobenius} has stated the following theorem in the memoir mentioned:

VI. \gesperrt{Every group of degree $p^{\alpha} q$, if $p$ and $q$ are primes different from one another and $\alpha$ denotes an arbitrary positive exponent, is metacyclic.}

The following proof of this theorem comes from a communication by letter from \gesperrt{Frobenius} to the author.

First of all it is clear that it suffices to prove that no group $P$ of degree $p^{\alpha} q$ is simple. For our theorem is proved for $\alpha = 1$. If we therefore assume it already proved for every smaller $\alpha$, and suppose that $P$ has a proper normal divisor $Q$ which comprises more than the unit element, then the two groups $Q$, $P/Q$ are metacyclic by the assumption or by Theorem IV, and hence $P$ is also metacyclic.

We now assume, then, that the group $P$ of degree $p^{\alpha} q$ is simple, and we have to derive a contradiction from this assumption in order to establish its impossibility.

$\alpha)$ First, by Theorem I., \S.\,29, we have a divisor of $P$ of degree $q$. We will assume that there is a
%CONT>>
